% !TEX root = ../main.tex
\section{Conclusiones}
\label{sec:conclusiones}

El objetivo principal de este trabajo es el desarrollo de modelos de inteligencia artificial capaces de predecir el carbono que una parcela forestal capturará en un horizonte temporal determinado, a partir de información estructural, edáfica, climática y espectral recogida en los Inventarios Forestales Nacionales y en fuentes de datos auxiliares. Los resultados obtenidos permiten afirmar que dicho objetivo se ha alcanzado de manera satisfactoria.

\subsection*{Selección de variables y preprocesamiento}

El proceso de selección manual de variables (sección~\ref{sec:entrenamiento}) permitió reducir un conjunto inicial de 445 predictores candidatos a un subconjunto compacto de 44 variables (véase Tabla~\ref{tab:entradamodelo}), manteniendo una representación equilibrada de todos los ámbitos ecológicos implicados: estructura de masa, composición de especie, propiedades edáficas, terreno, climatología e índices de vegetación espectral. Esta reducción no solo facilitó el entrenamiento de modelos más eficientes, sino que garantizó la interpretabilidad ecológica de los predictores seleccionados.

\subsection*{Rendimiento de los modelos}

Los modelos entrenados con los datos del IFN3 como variables explicativas obtienen, en general, las mejores métricas (véanse Tablas~\ref{tab:base_comp_ifn3_c4} y~\ref{tab:base_comp_ifn3_carbono_bruto4}). Esto se explica por la mayor calidad, homogeneidad metodológica y riqueza de variables recogidas en este inventario frente al IFN2. Para predicciones a medio plazo (entre 5 y 17 años), los modelos alcanzan niveles de precisión compatibles con aplicaciones prácticas de planificación forestal.

En cuanto a la variable objetivo, la predicción del carbono en toneladas absolutas (\texttt{carbono\_bruto4}) resulta sistemáticamente más precisa en términos de métricas globales que la predicción del carbono por unidad de superficie (\texttt{c4}, tC/ha), como se puede observar comparando las Tablas~\ref{tab:base_global_c4} y~\ref{tab:base_global_carbono_bruto4}. Sin embargo, el análisis por deciles (Figuras~\ref{fig:SMAPE_RMSE_deciles_c4_carbono_bruto4} y~\ref{fig:distribucion_carbono_bruto4_combinado}) revela que para instancias con valores de carbono muy bajos (primera biomasa), los errores relativos (SMAPE) son elevados en ambas variables con independencia del modelo empleado. Los modelos ofrecen sus mejores resultados en el rango intermedio-alto de biomasa, donde el balance entre cantidad de datos, variabilidad intrínseca y magnitud de la variable objetivo es más favorable.

El análisis de la variable \texttt{periodo} (Figuras~\ref{fig:LightGBM_rmse_y_pct_rmse_cuantiles}, \ref{fig:Stack5_MLP_rmse_y_pct_rmse_cuantiles}, \ref{fig:CatBoost_combined} y~\ref{fig:Stack5_MLP_combined_rmse_y_pct_rmse_cuantiles}) muestra que el error crece con el horizonte temporal de predicción, especialmente cuando los datos de entrenamiento proceden del IFN2 (periodos superiores a 17 años). No obstante, los errores medianos (cuantil 50) se mantienen relativamente estables a lo largo de todos los periodos, lo que indica que los modelos asimilan correctamente el comportamiento medio del crecimiento forestal incluso para horizontes largos. Los errores extremos (cuantil 90) son los más sensibles a la calidad de los datos y al alejamiento temporal.

El análisis de sensibilidad sobre escenarios representativos (Figuras~\ref{fig:sensibilidad_escenarios} y~\ref{fig:sensibilidad_escenarios_carbono_bruto4}) confirma que los modelos capturan de forma cualitativa correcta la dinámica de crecimiento: la tendencia general es ascendente con el periodo, y los modelos reproducen los efectos de los cambios de inventario (IFN2 a IFN3) observados en los datos reales.

\subsection*{Modelos globales frente a modelos locales}

La comparación entre modelos ``globales'' (entrenados con IFN2 e IFN3 conjuntamente) y ``locales'' (entrenados con un único inventario), mostrada en las Figuras~\ref{fig:comparativa_mejores}--\ref{fig:disc_sub4}, revela un comportamiento diferenciado según el inventario de validación. Para el IFN3, los modelos locales superan a los globales, ya que la inclusión de datos del IFN2 introduce heterogeneidad que penaliza la predicción sobre datos de mayor calidad. Para el IFN2, en cambio, los modelos globales son superiores, pues se benefician de los patrones adicionales aprendidos con los datos del IFN3. Este resultado subraya que la elección del conjunto de entrenamiento debe adaptarse al horizonte temporal y al inventario de aplicación.

\subsection*{Modelos de \textit{stacking}}

La incorporación de esquemas de \textit{stacking} (Tablas~\ref{tab:stack_configs} y~\ref{tab:meta_modelos}) produce mejoras sistemáticas pero modestas respecto a los mejores modelos base (véanse Tablas~\ref{tab:stack_global_c4} y~\ref{tab:stack_global_carbono_bruto4} frente a Tablas~\ref{tab:base_global_c4} y~\ref{tab:base_global_carbono_bruto4}). El conjunto de stacking con configuración 5 y metamodelo MLP es, de forma consistente, el que mejores métricas obtiene tanto para \texttt{c4} como para \texttt{carbono\_bruto4}. El hecho de que la mejora sea sistemática en todos los casos descarta que se trate de un artefacto estadístico, pero su magnitud limitada plantea la pregunta de si la complejidad adicional del stacking está justificada frente al uso del mejor modelo individual en un hipotético sistema en producción. Los algoritmos basados en \textit{gradient boosting}, especialmente CatBoost y LightGBM, son los que ofrecen el mejor rendimiento como modelos base en todos los escenarios evaluados.

\subsection*{Interpretabilidad}

El análisis SHAP (Figuras~\ref{fig:beeswarm_grouped_LightGBM} y~\ref{fig:beeswarm_grouped_CatBoost}) revela que la variable más influyente en las predicciones es la densidad arbórea (\texttt{npies}), seguida de la identidad de especie (\texttt{especie\_id}). Las variables climáticas y los índices de vegetación actúan como moduladores secundarios, refinando la predicción especialmente a escala estacional. Este resultado es coherente con la ecología forestal: la densidad y la especie son los determinantes primarios de la acumulación de biomasa, mientras que el clima y el estado espectral de la vegetación condicionan su tasa de crecimiento.

\subsection*{Comparación con la literatura}

En trabajos similares, como el de \citet{fasihi2024assessing} para el norte de Italia, se predicen valores de carbono forestal a partir de observaciones de teledetección sin uso de datos de campo medidos in situ. Aunque el contexto es distinto ---ellos predicen el carbono presente en el momento de la medición, sin proyección temporal--- los errores que reportan son notablemente superiores a los obtenidos en este trabajo, especialmente en el rango de predicción más corto (alrededor de 5 años), que es el más comparable. Esto pone de manifiesto el valor añadido de incorporar información estructural de campo (IFN) frente a enfoques basados exclusivamente en teledetección.

\subsection*{Limitaciones}

La principal limitación del estudio reside en la heterogeneidad e irregularidad de los datos disponibles. La distribución de instancias en función de la variable \texttt{periodo} es muy desigual: algunos años cuentan con miles de muestras mientras que otros tienen apenas unas pocas decenas. Esta irregularidad, inherente a la naturaleza de los Inventarios Forestales (levantados en periodos de tiempo distintos y con metodologías que evolucionan entre iteraciones), limita la capacidad del modelo para ciertos horizontes temporales, sobre todo en el rango de años con menos datos y donde la fuente principal es el IFN2. Pese a ello, los modelos demuestran capacidad de generalización razonable incluso en estas condiciones.

\subsection*{Trabajo futuro y aplicabilidad}

De cara a una posible aplicación en producción como sistema de apoyo a la recomendación de proyectos de carbono, sería necesario evaluar en qué medida las parcelas del IFN son representativas de plantaciones forestales destinadas a créditos de carbono, que pueden diferir en densidad, manejo y composición. Una extensión natural del modelo sería incorporar datos LiDAR, cuya capacidad para caracterizar la estructura tridimensional del dosel forestal ha demostrado en otros estudios mejorar significativamente las predicciones de biomasa y carbono \citep{nyakudanga_treemes}. Igualmente, la adaptación del modelo a otros contextos geográficos o su integración con datos del IFN4 como variable objetivo constituyen líneas de trabajo prometedoras.
